\section*{Open Questions}

\begin{description}
  \item[Q1] The pyqsp \texttt{PolySign.generate(chebyshev\_basis=True)}
        wrapper is broken in the current release (a local variable
        \texttt{scale} is referenced before assignment when
        \texttt{return\_scale=False}) -- we had to bypass it and build
        the Chebyshev interpolant of $\mathrm{erf}(\kappa x)$ manually
        via \texttt{PolyTaylorSeries}. How many downstream QSVT
        replication attempts silently produce garbage when they hit
        this exact code path? A cross-repo bisect on pyqsp API
        deprecations plus a formal ``QSP convention-mismatch smoke
        test'' would be worth publishing.
  \item[Q2] The ``\texttt{Wx}''-convention QSP+QSVT signal operator
        $[[x, i\sqrt{1-x^2}],[i\sqrt{1-x^2},x]]$ deposits an odd-degree
        real target polynomial in $\mathrm{Im}[U_\Phi[0,0]]$, not in
        $\mathrm{Re}[U_\Phi[0,0]]$. Many QSVT tutorials and even the
        pyqsp response module contain a comment saying ``sym\_qsp real
        component means little'' -- yet the widely-cited pedagogical
        rule of thumb is that ``$\mathrm{Re}[U[0,0]] = P(x)$''.
        Investigating exactly which QSVT papers/pedagogical treatments
        (including the paper under replication and its predecessors)
        adopt which convention -- and whether the confusion has caused
        undetected sign/parity errors in downstream numerical
        replications of Chebyshev-history-state or LCU-QSVT circuits --
        seems useful.
  \item[Q3] The paper's ``$\kappa_S$'' (basis condition number of the
        Jordan form) enters QEVT complexity multiplicatively, and the
        paper's convergence claims implicitly assume $\kappa_S$ is
        moderate.  How does the QEVT approach degrade empirically as
        $\kappa_S$ grows on a controlled family of non-normal test
        matrices (e.g.\ Jordan blocks perturbed by $\epsilon$, or the
        Grcar / Kahan pseudospectra classics)?  This is out of reach of
        a Hermitian-only single-day replication but is critical for
        practitioners trying to apply QEVT.
  \item[Q4] The pyqsp \texttt{sym\_qsp} Newton solver converges to
        residual $<10^{-12}$ in 5--8 iterations for our targets
        (degrees $\le 41$).  How does its scaling break down at
        production degrees ($\sim 10^3$--$10^4$) required for the
        paper's ODE-solver application with tight $\epsilon$?  Do any
        of the residual failure modes require the extended-precision
        or preconditioned solvers of Ying \& Kunis or Chao--Ni?  A
        head-to-head comparison of \texttt{sym\_qsp} vs.\ QSPPACK vs.\
        pennylane's \texttt{qml.qsvt} at high degree would help pin
        down the practical regime where the paper's asymptotic
        complexity kicks in numerically.
  \item[Q5] The paper's key subroutine -- Chebyshev history state
        generation via a matrix generating function -- is presented as
        a block-encoded circuit but not (as far as we saw) demonstrated
        on a small classically-simulable instance.  Building even a
        toy 2--3 qubit Chebyshev history state circuit (e.g.\ for
        $H = \sigma_x$ or a $4\times4$ Toeplitz-shift matrix) and
        verifying that the LCU-of-Chebyshev structure matches the
        paper's block-encoding recipe would give the community a
        canonical worked example.  Such an example is a natural next
        step to move beyond our Hermitian-QSVT-equivalent spot-check.
\end{description}
